% Cote 119: offprints, reports, courses, CV and letters (typed and printed part): archive cote 119, site page 283 (PDF page 284).
% typed letter. Draft compared: none.
% AI-transcribed by Claude Opus 5.5 (claude-opus-5-5, adaptive thinking, effort high), Message Batches API, batch r119_full (msgbatch_01DwZNm5Y7YsKBWyAN4sFs8w), 2026-11-13; not yet checked by eye.
\scanpage{119/283}
of order 2 .

One of the projects here at Bangor is the Popularisation of
Mathematics. There is a mixture of motives, of course, one of which is
advertisement for Bangor, and so to improve the Department and its work.
But the main motive is relevant to your writings in part of Recollte et
Semaille, and is to emphasise the notion of joy in Mathematics, rather than
pure competition, and to convey this to as wide a public as possible. I
enclose a copy of our Mathematics and Knots, and of a review of it. It took
us 4 years to get the exposition to this stage! I believe the hard work and
the underlying structure shows through.

I have been hearing from others of you, and there is continuing
interest in `Pursuing Stacks'. It may be that the idea of editing it is not
a sensible one. Baues suggested that it should be reproduced as it is,
complete with pen alterations, so as to give an authentic impression of the
work. I am not sure who would publish it, but it would be interesting to
know of your reactions to this idea. Would you like this to be
followed up?

There could be extra chapters of comments, guides, and so on. For
example, I have just had, as part of a long correspondence about cubical
sets and $\infty$-categories, a request for a copy from a computer scientist,
Vaughan Pratt, at Stanford. For me, this is suggestive that the idea of
`higher dimensional algebra', of a logic which works in 2 or more
dimensions, is an idea whose time has just about come.

In the UK, such an idea tends to be regarded as, how shall I say,
foreign, or French (!), since the English are most pragmatic, and think
Mathematics should develop from already stated and formulated problems. But
when the powers that be ask what are the strengths and weaknesses of UK
Mathematics, I ask in return in which country originated sheaves, spectral
sequences, fibre bundles, K-theory, $\lambda$-rings, schemes, bordism, jets,
foliations, and, to be more recent, cyclic cohomology, cat$^{n}$-groups (the last
being not so famous yet). About the only \textit{basic} technique at that level I
can think of from the UK is the ideas of Whitehead on combinatorial
homotopy, and algebraic \textit{K}-theory, and, I would like to believe for the
future, applications of multiple groupoids. I exclude here the Adams
spectral sequence, since this is an advanced technique in a specialist area,
rather than a concept opening up new areas. (The specialists might
disagree.) I wrote this and other comments in a letter to the Science and
Engineering Research Council, since they had asked generally for views, and
was thanked for being so frank, which might mean `some people say more than
is necessary, and don't help themselves'.

I sent in a Research Proposal of the general kind on `Multiple
groupoids and category theory', putting it at the highest level of a
programme, without too much detail on problems, basically because I would
prefer it to be rejected at that level, rather than a detailed problem set.
At least I have put my argument, and let history say which of us is right!
However, this was rejected. I guess I cannot get the tone in the way that
my colleagues expect. Also, there is a credibility gap: can the notion of a
generalisation of groups really exist? Can there be a concept of higher
order symmetry, and higher order algebra, with wide applications? If it did
exist, it would be very important. But it is not well known. Therefore,...

There are problems in carrying out the somewhat vaguer, next stage of
the programme. The original ideas in homotopy theory have to a large extent
been encapsulated in the Van Kampen Theorem with Loday, and the question is
whether to concentrate on more detailed development, or to push the general
programme, and if so, in what way. So I am feeling my way, but with a large
number of immediate projects. Essentially, I have in view a clear
foreground, some distant peaks, and a mushy middle ground. Do I cultivate

\mspage{2}

\margin{283}
